\section{First-principles derivation candidate: one-loop self-energy}
\label{sec:eq-self-energy}

Source:
\texttt{Mathematica\_Notebooks/Quantum\_Mechanics/r\_e\_derivation\_self\_energy.wl}
(179 lines, scaffold for issue~\#64). Companion to
\texttt{Equation\_Verification/Dual\_Relativistic\_Quantum\_Mechanics\_I.md}~\S III.D.

\paragraph{Target.}
$r_e/r_0 = 0.4994205099128317 \pm 2.5\times 10^{-13}$ (PR~\#62 triangulation).

\paragraph{Cell~1 --- standard QED one-loop on-shell mass shift.}
Using the Schwinger proper-time representation,
$i/(A - i\epsilon) = \int_0^\infty ds\, e^{-isA}$, the on-shell mass shift is
\[
\frac{\delta m}{m} \;=\; \frac{3 \alpha}{4\pi}
   \left[\,\ln\!\frac{\Lambda^2}{m^2} \;+\; \frac{1}{2}\,\right]
\qquad \text{(Bjorken--Drell Eq.~10.59).}
\]
At $\alpha = 1/137.036$, $\Lambda/m = 1$ (Bethe NR cutoff):
$\delta m / m \approx 8.71\times 10^{-4}$.
\wolframcheck{Confirmed 2026-05-26.}

\paragraph{Cell~2 --- dual framework, one-loop, closed-form inversion.}
The DRQM~I (III.22) $g$-factor formula reads
$g_r(x) = 2\bigl[1 - 4/(2x + 1)\bigr]$.
Setting $g_r(x) = -2(1 + a_e)$ with
$a_e \to a_e^{(1)} = \alpha/(2\pi)$ (Schwinger 1948) and solving for $x$:
\[
\boxed{\;
  \left.\frac{r_e}{r_0}\right|_{1\text{-loop}}
  \;=\; \frac{2 - \alpha/(2\pi)}{4 + \alpha/\pi}
  \;=\; \frac{2 - a_e^{(1)}}{2(2 + a_e^{(1)})}
\;}
\]
\wolframcheck{\texttt{FullSimplify} on $r_e/r_0 - (2 - a_e^{(1)})/(2(2 + a_e^{(1)}))$
returns $0$.}
At $\alpha = 1/137.036$, 18-digit: $0.499419632159865\,6$.

\paragraph{Cell~3 --- comparison to triangulated target.}
\begin{center}
\begin{tabular}{lll}
\toprule
Order               & $r_e/r_0$               & residual vs.~triangulated \\
\midrule
Dirac tree ($a_e=0$) & $0.5$                  & $+5.79\times 10^{-4}$ \\
1-loop Schwinger    & $0.499419632159865\,6$  & $-8.78\times 10^{-7}$ \\
2-loop Sommerfeld   & $0.499420517282276\,9$  & $+7.37\times 10^{-9}$ \\
\bottomrule
\end{tabular}
\end{center}
2-loop is closer than 1-loop by a factor $\approx 120$.
\wolframcheck{Confirmed 2026-05-26.}

\paragraph{Cell~4 --- higher-loop convergence.}
Including $a_e^{(2)} = -0.328478965579193\,(\alpha/\pi)^2$,
$a_e^{(3)} = +1.181241456587\,(\alpha/\pi)^3$,
$a_e^{(4)} = -1.91245\,(\alpha/\pi)^4$, and
CODATA-full $a_e^{\rm expt} = 0.001\,159\,652\,18\,059$
into the closed-form cutoff
$r_e/r_0 = (2 - a_e)/(2(2 + a_e))$
drives the residual against the triangulated target below
$10^{-12}$ by 4-loop; CODATA-full reproduces the triangulated value to its full
precision (precision floor $\sigma = 2.5\times 10^{-13}$).

\paragraph{Hypothesis split (the load-bearing question for Tepper).}
\begin{itemize}[topsep=2pt,itemsep=0pt]
  \item \textbf{Hypothesis (ii)} --- the dual framework's proper-time one-loop
        vertex correction \emph{reproduces} the standard-QED Schwinger anomaly
        identically (photon propagator unchanged; dual structure absorbed into
        the matter-side (II.3) ``potential-in-the-mass'' kernel which at
        one-loop is formulation-independent). Under this hypothesis, the
        closed-form $r_e/r_0 = (2 - a_e)/(2(2 + a_e))$ matches the
        triangulated value once $a_e$ is supplied by CODATA. This is the
        \emph{reproduction-by-inheritance} reading.
  \item \textbf{Hypothesis (i)} --- the dual framework's proper-time one-loop
        vertex (with the dual photon propagator and bound-state structure)
        is computed \emph{independently} and shown to reproduce
        $a_e^{(1)} = \alpha/(2\pi)$. This unconditional reading is tracked in
        issue~\#75.
\end{itemize}

\paragraph{Crocco compliance.}
Cells~1--4 arithmetic is pragmatic AI. The Cell~2 dual-framework hypothesis
identification (hypothesis (ii)) is substantive AI; the Tepper-blocker is
which of (i) or (ii) is the framework's intended position. This is the
principal question carried into the closing list.

\authorreview{The Cell~2 derivation under hypothesis (ii) reproduces the
Schwinger closed-form; does the dual framework's intended position match
hypothesis (ii), or does it require hypothesis (i) (an independent
proper-time vertex computation)? Issue~\#75 tracks the (i) path.}
